% ============================================================
% АННОТАЦИЯ (заглавный лист — большая рамка Форма 2)
% ============================================================
\thispagestyle{gostfirst}
% Дополнительный нижний отступ под штамп Формы 2 (40 мм вместо 15 мм)
\enlargethispage{-25mm}

\section*{АННОТАЦИЯ}
\addcontentsline{toc}{section}{АННОТАЦИЯ}

Пояснительная записка \pageref{LastPage}~стр., 12~рис., 8~табл., 1~приложение, 50~использованных источников.

Работа посвящена разработке интерактивной обучающей веб-платформы SopromatLab для изучения основ сопротивления материалов с визуализацией напряжённо-деформированного состояния статически определимых балочных конструкций. Формализованы математические модели трёх типов балок, разработаны алгоритмы расчёта эпюр $Q(x)$, $M(x)$, $N(x)$, $\theta(x)$, $v(x)$ методом сечений и численного интегрирования; реализована автоматическая верификация по дифференциальным зависимостям Журавского. Платформа включает интерактивный редактор, теорию, тренажёр, учебные сценарии, пошаговое решение, контекстный AI-консультант на базе большой языковой модели (LLM), экспорт в~PDF и обмен по~ссылке.

На защиту выносится графическая часть в~объёме 7~листов формата~А1.

\textbf{Ключевые слова:} сопротивление материалов, напряжённо-деформированное состояние, эпюры внутренних усилий, интерактивная обучающая платформа, метод сечений, зависимости Журавского.

\newpage
